\section{Background and Related Work}
\label{sec:background}

%% Budget: ~0.9 pp. Three threads, each ending in the one sentence that separates this work from
%% it. Every \cite key is in references.bib with a `% verified:` record (sweep 2026-06-12,
%% re-verified 2026-08-30; six entries added 2026-09-09). 2026-09-13 (plan §13): tightened twice,
%% no citation dropped. 2026-09-13 readability pass: sentences split, no citation changed.

\subsection{Erosion, conformance, and architecture as code}

Perry and Wolf distinguished architectural \emph{erosion}, violations of the intended architecture,
from \emph{drift}, the loss of coherence that comes from insensitivity to
it~\cite{perry1992foundations}. Surveys~\cite{desilva2012erosion,li2022erosion} and practitioner
studies~\cite{anthony2024drift,ali2018consistency} still name outdated documentation as a main
cause and reliable documentation as the main remedy. The classical counter-measures compare an
intended model against the code: reflexion models~\cite{murphy2001reflexion},
dependency-structure matrices~\cite{sangal2005dsm}, static conformance
checking~\cite{knodel2007comparison,passos2010conformance}, lately run continuously in the
delivery pipeline~\cite{bucaioni2024conformance}, and commercial analysers that check rules over
the full dependency graph~\cite{understand,ndepend}. All presuppose the intended model as input,
which is exactly what large legacy systems lack. The C4 model~\cite{brown2026c4} and the
Structurizr DSL~\cite{structurizrdsl} made textual, version-controlled \emph{architecture as
code}~\cite{bucaioni2025aac} mainstream. \tool{} produces that artifact from evidence rather than
by hand, to regenerate byte-identically and drift-check rather than to draw once.

\subsection{Software architecture recovery}

The dominant automatic recovery line is bottom-up clustering of a dependency
graph~\cite{ducasse2009sar,shtern2012clustering,sarhan2022clustering}. Its families include Bunch
(search for maximal modularisation quality, MQ)~\cite{mancoridis1998mq,bunch}, ACDC
(comprehension-driven subgraph patterns)~\cite{acdc}, WCA (agglomerative)~\cite{wca}, LIMBO
(information-theoretic)~\cite{limbo}, FCA (a fast matrix
heuristic)~\cite{teymourian2022fca,zhang2023sarif}, and ARC (structure plus LDA topics over
identifiers and comments)~\cite{arc}. We run ACDC, WCA, LIMBO, and ARC as baselines
(\S\ref{sec:design}). We do not run Bunch, which \textsc{Arcade} does not bundle and the recent
evaluations do not compare against; its MQ objective appears only as a descriptive statistic of
the build graph (\S\ref{sec:rq1}). \textsc{Arcade}~\cite{arcade} bundles the others with the
standard metrics: MoJoFM, the normalised move-and-join distance between
partitions~\cite{tzerpos1999mojo,mojofm}, c2c, the per-cluster overlap~\cite{garcia-comparative},
and a2a, which additionally charges for adding, removing, and moving clusters~\cite{le2015a2a}.
The metrics disagree by design: MoJoFM favours many small clusters because merging is cheap,
whereas a2a varies little~\cite{lutellier2018measuring,zhang2023sarif}. Zhang et al.\ therefore
brought the chance-corrected adjusted Rand index (ARI), which charges over-splitting in full, into
SAR evaluation~\cite{zhang2023sarif,jabeen2026gaer}.

Ground truth is the field's bottleneck: certified references take architect-weeks to
obtain~\cite{garcia2013groundtruth}, the public set is a handful of C/C++ and Java
systems~\cite{lutellier-deps,zhang2023sarif}, and none is C\#. Lutellier et al.\ showed that the
accuracy of the \emph{dependencies} fed to a technique matters as much as the
technique~\cite{lutellier-deps,lutellier2018measuring}, a finding our build-graph-first design
takes to its conclusion. Stability is the second gap: Wu et al.\ found the most stable algorithms
the least authoritative and vice versa~\cite{wu2005comparison}, and ARC moves by tens of MoJoFM
points between identical inputs even after seeding~\cite{link2019value}. Our RQ4 makes
byte-identical regeneration a tested invariant rather than a patch.

The current frontier fuses further sources into the structural graph: multiple dependency
models~\cite{altinisik2024multidep}, code text and folders (SARIF)~\cite{zhang2023sarif},
pattern knowledge (SemArc)~\cite{zhao2026semarc}, code models and graph auto-encoders (SSAR,
REFINE, GAER)~\cite{ding2026ssar,zhang2026refine,jabeen2026gaer}, or class-name embeddings
(ClassLAR)~\cite{he2025classlar}; all treat structure as unknown and infer it from files.

Build systems have been reverse-engineered before. Tu and
Godfrey's build-time architectural view~\cite{tu2001buildtime} and Adams et al.'s MAKAO recover
the build dependency graph and its evolution~\cite{adams2007makao,adams2008linuxbuild}. The build
maintenance line quantifies how much that graph
churns~\cite{mcintosh2011buildeffort,mcintosh2012javabuild,nejati2024buildchanges}, derives
change-impact graphs from CMake and Make~\cite{meidani2023dipidi}, and finds declared build
dependencies error-prone~\cite{lyu2024builddeps}, one reason we refine them with source evidence.
None turns the graph into an architecture model, curates it, or checks it for drift. The nearest recent work, the Repository Intelligence Graph, reads the CMake File API to
give coding agents a map of a repository~\cite{chernyshahar2026rig}. It too produces no
architecture model, has no curation or drift layer, covers neither MSBuild nor C\#, and is not
evaluated against reference architectures.

\subsection{LLMs for architecture recovery and documentation}

For recovery, ArchAgent lets agents synthesise multi-view diagrams of large repositories, judged
by expert F1 rather than partition metrics; the tool is not released, and its own ablation reports
that dependency context improves accuracy~\cite{archagent}. SemRef post-processes the output of
ten SAR techniques with an LLM and raises MoJoFM substantially~\cite{zhang2026semref}, so the
model edits structure after the fact. CIAO~\cite{deluca2026ciao} and CodeWiki~\cite{codewiki2026}
generate architecture documentation end-to-end from code. Hybrids on C++~\cite{hatahet2025sad},
ROS~2~\cite{benchat2026ros2}, and industrial Java~\cite{rukmono2026dsar} run deterministic
extraction first and let the LLM select or abstract the elements. The evidence on what LLMs do
well is consistent: they identify styles but struggle with
relationships~\cite{amalfitano2025recovery}, 4{,}137 generated views over 340 repositories sat at
code level rather than architectural abstractions~\cite{sathvika2026views}, and a reproduction
attempt of 18 LLM-based software-engineering studies reproduced none in
full~\cite{angermeir2026reproducibility}.

\tool{} takes the opposite position on the one point every LLM-recovery line shares: structure
comes deterministically from the build graph and refining source evidence, and the model may only
name and describe what that evidence already contains (\S\ref{sec:approach}). The nearest
precedent for that division of labour is a recent drift detector that diffs runtime traces
deterministically and confines the LLM to writing the report~\cite{deangelis2026drift}. We
evaluate with the community's own partition metrics and ground truths and add what the
LLM-centric line does not report: reproducibility under rerun, stability under code change, and
drift-detection accuracy.
